\section{数据预处理与统计特征}

\subsection{数据完整性检查}

读取Excel文件后，以$(d,h)$作为联合主键检查记录数量、日期范围和小时范围。附件实际包含30个日期、每个日期24个小时且无重复键。对每个日期计算
\begin{equation}
A_d=\sum_{h=0}^{23}a_{dh},\qquad \bar a_d=\frac{A_d}{24},\qquad a_d^{\max}=\max_h a_{dh}.
\end{equation}
全月总量为2155604件，日总量均值为71853.47件，标准差由程序从原始数据直接计算。所有数据均保持件为基本计量单位，不进行未经说明的归一化。

\subsection{日内波动与峰值}

从热力图可以观察到小时6附近存在低谷，而0--4时和20--23时多次出现高值。日总量与峰值并不完全同步：第12天总量107377件，同时也是两种排班模型的峰值日；第26天日总量最低，为59156件。排班模型因此必须同时响应日总量和到货时序。

\subsection{下界与模型合理性}

问题一中每位工人每班最多处理200件，给出日人数下界$\lceil A_d/200\rceil$；问题二中每位工人每班处理$7\times25+10=185$件，给出$\lceil A_d/185\rceil$。这些下界只考虑全天总量，不考虑班次覆盖和16时截止。将求解结果与下界逐日对比，是检验整数规划结果是否被过度约束的重要方法。问题二的增量主要集中在需要早班产能的日期，说明增量具有结构性而非随机误差。

\begin{figure}[htbp]
\centering
\includegraphics[width=.88\textwidth]{figures/shift_start_frequency.pdf}
\caption{问题二班次起点在30天中的出现频率}
\label{fig:shiftfreq}
\end{figure}

班次起点频率表明16时班次在30天中始终被选择，这是因为23时仍有货物需要完成；其余班次起点根据当日早晚峰值动态变化。这一现象也解释了固定五班起点方案相对动态排班的损失。
